% Generated; do not edit.
\setresult{embedding_caption/plantnet/species}{34}
\setresult{embedding_caption/plantnet/rare}{11}
\setresult{embedding_caption/plantnet/without_photo}{1}
\setresult{embedding_caption/plantnet/cap}{30}
\setresult{embedding_caption/plantnet/min_images}{4}
\setresult{embedding_caption/plantnet/text}{Embedding structure of the 34 species of the family \textit{Ranunculaceae} in the Pl@ntNet300K dataset: up to 30 test images per species and the species' class means, at the angle to the family's spherical median (rings every 15 degrees). The radial axis of each panel runs from 50 to 75 degrees (Flat), 45 to 75 degrees (Hierarchical), 45 to 70 degrees (Independent) and 5 to 45 degrees (Conditional); images beyond the outer limit are placed on it. Hollow markers (faded in the photo figures) show the 11 species with fewer than 4 images. No photo is available for 1 species.}
\setresult{embedding_caption/global_lepi/species}{68}
\setresult{embedding_caption/global_lepi/rare}{0}
\setresult{embedding_caption/global_lepi/without_photo}{2}
\setresult{embedding_caption/global_lepi/cap}{30}
\setresult{embedding_caption/global_lepi/min_images}{4}
\setresult{embedding_caption/global_lepi/text}{Embedding structure of the 68 species of the family \textit{Pterophoridae} in the Global Lepidoptera dataset: up to 30 test images per species and the species' class means, at the angle to the family's spherical median (rings every 15 degrees). The radial axis of each panel runs from 30 to 90 degrees (Flat), 30 to 90 degrees (Hierarchical), 35 to 80 degrees (Independent) and 30 to 75 degrees (Conditional); images beyond the outer limit are placed on it. No photo is available for 2 species.}
